\section{Objetivos}\label{sec:objectives}

O objetivo central da pesquisa é criar e avaliar modelos de inteligência artificial que identifiquem automaticamente as bolhas em imagens do céu noturno do programa EMBRACE do INPE, aprimorando o monitoramento do clima espacial brasileiro.

Os passos práticos incluem:
\begin{enumerate}
   \item Organizar um conjunto de imagens do INPE, contendo noites com e sem bolhas sob diferentes condições climáticas.
   \item Preparar as imagens removendo luzes indesejadas e destacando as estruturas de interesse.
   \item Diminuir o tamanho dos arquivos de imagem matematicamente, mantendo o formato e a localização espacial das informações.
   \item Treinar um comitê de inteligência artificial que une três algoritmos diferentes, utilizando estratégias para equilibrar os dados de noites normais e noites com bolhas.
   \item Validar as decisões da máquina para confirmar se ela avalia o formato correto do fenômeno e testar sua estabilidade durante monitoramentos contínuos de longos períodos.
\end{enumerate}